\subsection{Flexible and Mixture Discriminant Analysis}

Linear discriminant analysis is useful partly because it is simple: each class is represented by a centroid, distances are measured in a common covariance metric, and the decision boundaries are linear.  Its limitations are equally structural.  A single centroid may not represent an inhomogeneous class, linear boundaries may be too rigid, and high-dimensional predictors may require smoothness or shrinkage beyond ordinary covariance pooling.  Flexible, penalized, and mixture discriminant methods keep the LDA comparison geometry while changing the regression fit, the metric, or the number of class prototypes.

\begin{method} \label{meth:sl-flexible-penalized-discriminant-analysis} \hlt{Flexible and Penalized Discriminant Analysis}\\
Optimal scoring represents LDA by regressions from $x_i$ to transformed labels.\\
Let $g_i\in\{1,\ldots,K\}$, $1\leq L\leq K-1$, label scores $\theta_\ell:\{1,\ldots,K\}\to\R$, and fitted score maps $\eta_\ell:\mathcal{X}\to\R$, with centering and orthogonality constraints on the $\theta_\ell$ removing constant score directions.
\begin{enumerate}[label=(\roman*)]
\setlength{\itemsep}{0pt}
\item \hlt{Flexible discriminant analysis}: minimize $\mathcal{L}_{\mathrm{FDA}}=n^{-1}\sum\limits_{\ell=1}^{L}\sum\limits_{i=1}^{n}\{\theta_\ell(g_i)-\eta_\ell(x_i)\}^2+\lambda\sum\limits_{\ell=1}^{L}J(\eta_\ell)$ over $\theta_\ell$ and $\eta_\ell$ from a chosen smoother class: linear, spline, additive, MARS, kernel, or similar.
\item \hlt{Penalized discriminant analysis}: impose $\eta_\ell(x)=h(x)^\top\beta_\ell$ and roughness penalty $\beta_\ell^\top\Omega\beta_\ell$.\\
Equivalently classify by $D_\lambda(x,k)=\bigl(h(x)-\bar{h}_k\bigr)^\top(\widehat{\Sigma}_{W,h}+\lambda\Omega)^{-1}\bigl(h(x)-\bar{h}_k\bigr)$, with $\bar{h}_k$ and $\widehat{\Sigma}_{W,h}$ the class-$k$ centroid and pooled within-class covariance in $h$-coordinates.
\end{enumerate}
For $\eta_\ell(x)=x^\top\beta_\ell$ and no roughness penalty, the fitted directions recover LDA up to scale; FDA changes the regression class to obtain nonlinear boundaries, while PDA shrinks rough or unstable $h$-directions.
\end{method}

\begin{breakablealgorithm}
\caption{FDA Coordinates by Optimal Scoring}
\label{alg:sl-fda-optimal-scoring}
\begin{algorithmic}
\Require Class labels $g_i\in\{1,\ldots,K\}$, indicator response matrix $Y\in\R^{n\times K}$, linear smoother $S_\lambda$ for the chosen regression method
\Ensure Nonconstant discriminant coordinate functions $\eta_1,\ldots,\eta_{K-1}$
\State Fit the multivariate regression of $Y$ on the features: $\widehat{Y}\gets S_\lambda Y$
\State Let $\eta^*(x)$ denote the $K$ fitted regression functions evaluated at a new input $x$
\State Set $D_\pi\gets Y^TY/n$
\State Solve $Y^T\widehat{Y}\theta=\gamma D_\pi\theta$ and normalize eigenvectors so that $\Theta^TD_\pi\Theta=I$
\State Discard the constant scoring direction
\For{each remaining eigenvector $\theta_\ell$}
\State Set $\eta_\ell(x)\gets \theta_\ell^T\eta^*(x)$
\EndFor
\end{algorithmic}
\end{breakablealgorithm}

\begin{remark} \label{rem:sl-fda-basis-expansion-lda-link} \hlt{FDA as LDA After a Regression Map}\\
Let $\eta=(\eta_1,\ldots,\eta_L):\mathcal{X}\to\R^L$ and $z=\eta(x)$.  FDA applies LDA in $z$-space, so linear discriminants in $z$ pull back to decision boundaries in $x$ determined by the regression class for $\eta$.\\
Polynomial $\eta$ gives algebraic boundaries, quadratic $\eta$ gives quadratic boundaries, and additive, spline, or kernel $\eta$ gives the corresponding structured nonlinear boundaries.
\end{remark}

\begin{method} \label{meth:sl-mixture-discriminant-analysis} \hlt{Mixture Discriminant Analysis}\\
Mixture discriminant analysis replaces the class-$k$ LDA prototype by $R_k$ Gaussian subclasses while sharing one covariance $\Sigma$.  Let $\pi_{kr}\geq0$, $\sum_{r=1}^{R_k}\pi_{kr}=1$, centers $\mu_{kr}$, class prior $\rho_k$, $f_k(x)=\sum_{r=1}^{R_k}\pi_{kr}\phi(x;\mu_{kr},\Sigma)$, and $p_k(x)=\Pr(G=k\mid X=x)=\rho_k f_k(x)/\sum_{j=1}^{K}\rho_j f_j(x)$. \\
If $R_k=1$ for all $k$, this is ordinary common-covariance Gaussian discrimination.
\begin{enumerate}[label=(\roman*)]
\setlength{\itemsep}{0pt}
\item \hlt{E-step}: for $g_i=k$, set $r_{ikr}=\pi_{kr}\phi(x_i;\mu_{kr},\Sigma)/\sum_{s=1}^{R_k}\pi_{ks}\phi(x_i;\mu_{ks},\Sigma)$ for $r=1,\ldots,R_k$.
\item \hlt{M-step}: update $\pi_{kr}$, $\mu_{kr}$, and $\Sigma$ by weighted Gaussian maximum likelihood with weights $r_{ikr}$.
\item \hlt{Blurred-response variant}: use subclass columns $c_{kr}$ and row $i$ equal to $(r_{ik1},\ldots,r_{ikR_k})$ on class-$k$ columns, zero elsewhere; then fit weighted DA with optional rank restrictions, FDA, or PDA.
\end{enumerate}
\end{method}

\begin{remark} \label{rem:sl-mda-versus-prototype-rules} \hlt{MDA and Prototype Rules}\\
MDA is a probabilistic prototype classifier with prototypes $\mu_{kr}$, weights $\pi_{kr}$, and shared metric $\Sigma^{-1}$. It assigns $x$ by posterior scores $p_k(x)\propto \rho_k\sum_{r=1}^{R_k}\pi_{kr}\phi(x;\mu_{kr},\Sigma)$, not by a hard nearest-center rule.\\
Relative to classwise $k$-means, it estimates $\pi_{kr}$ and $\Sigma$; relative to an unconstrained Gaussian mixture, each component $(k,r)$ is nested inside the observed class $k$.
\end{remark}

\begin{remark} \label{rem:sl-discriminant-methods-computation} \hlt{Computational Scale}\\
Let $N$ be the sample size, $p$ the predictor dimension, and $T_{\mathrm{fit}}$ the cost of the chosen smoother.
\begin{enumerate}[label=(\roman*)]
\setlength{\itemsep}{0pt}
\item \hlt{LDA/PDA}: covariance-scale linear algebra costs about $Np^2+p^3$ before low-rank reduction.
\item \hlt{FDA}: cost is $T_{\mathrm{fit}}$; additive or adaptive basis fits can be nearly linear in $N$ for fixed model size, while dense kernel or spline smoothers can cost at least an $N\times N$ solve.
\item \hlt{MDA}: cost multiplies the common-covariance discriminant fit by EM iterations and depends on initialization because mixture likelihoods are nonconvex.
\end{enumerate}
\end{remark}
